% 本文件由 LaTeX 排版 Agent 根据有效 translation.json 自动生成。
% author=Apocrypha; work_id=1001; title=伯多禄福音

\chapter{\WorkTitle{伯多禄福音}}

% structure: p0001 source=h1 target=omitted reason=page-title-already-rendered-by-parent

1. 但犹太人中没有一人洗手，无论是\Person{黑落德}还是他的任何一位法官。当他们拒绝洗手时，\Person{比拉多}就站了起来。然后\Person{黑落德}王命令将主带走，对他们说：\Quote{凡我吩咐你们对他做的事，你们就去做。}\par

2. 那里站着\Person{若瑟}，他是\Person{比拉多}和主的朋友；他知道他们将要把他钉死，便来到\Person{比拉多}面前，请求领取主的遗体埋葬。\Person{比拉多}派人到\Person{黑落德}那里，请求领取他的遗体。\Person{黑落德}说：\Quote{\Person{比拉多}兄弟，即使没有人来求他，我们也打算埋葬他，尤其因为安息日临近了；因为法律上写着：\InnerQuote{太阳不落于被处死的人}。}\par

3. 他在无酵饼节的前一天，就是他们的庆节，把主交给了民众。他们抓住主，奔跑着推他，说：\Quote{我们既然获得了对他的权柄，就把他拖走吧。}他们给他穿上紫袍，让他坐在审判席上，说：\Quote{以色列王，请公正地审判。}有一个人带来一顶荆棘冠冕，放在主的头上。另一些人站着向他脸上吐唾沫，又有人打他的脸颊；还有人用芦苇刺他；有些人鞭打他，说：\Quote{让我们以此荣耀来尊崇天主子。}\par

4. 他们带来了两个凶犯，把主钉在他们中间。但他保持沉默，好像没有疼痛。当他们竖起十字架时，写了标题：\Quote{这是以色列王。}他们把他的衣服放在他面前，彼此分取，并为它们拈阄。其中一个凶犯责备他们说：\Quote{我们因所做的恶事而这样受苦，但这个人成了人类的救主，他对你们做了什么错事呢？}他们对他发怒，就下令不要打断他的腿，好让他在痛苦中死去。\par

5. 正午时分，黑暗笼罩了整个\Place{犹太地}；他们忧虑不安，生怕太阳在他还活着时落下；因为法律为他们写着：\Quote{太阳不落于被处死的人。}其中一个人说：\Quote{给他喝苦胆搀醋。}他们调和了，给他喝，就完成了所有事，并把他们的罪孽归到自己头上。许多人拿着灯走来走去，以为是黑夜，就跌倒了。主大声呼喊说：\Quote{我的能力，我的能力，你遗弃了我。}说完这话，他被接升。就在那时，\Place{耶路撒冷}圣殿的帐幔裂为两半。\par

6. 然后他们从主的双手拔出钉子，把他放在地上，整个大地震动，极大的恐惧临到。然后太阳照耀，发现是第九时辰；犹太人欢喜，把他的身体交给\Person{若瑟}去埋葬，因为他看到了他所行的善事。他领了主，洗净了，用细麻布裹好，带进自己的坟墓，那坟墓称为\Place{若瑟的花园}。\par

7. 那时，犹太人和长老以及司祭们，意识到他们对自己做了何等恶事，就开始哀叹说：\Quote{我们的罪有祸了！审判临近了，\Place{耶路撒冷}的结局也近了。}我和我的同伴们忧伤；我们心灵受伤，就躲藏起来；因为他们搜捕我们，如同凶犯，以为我们想要放火烧圣殿。做这一切事时，我们禁食，坐着悲痛，日夜哭泣，直到安息日。\par

8. 但是经师们、法利塞人和长老们彼此聚集，听到众人私下议论并且捶胸，说：\Quote{如果因他的死这些极大的征兆发生了，看他是多么正义}——长老们害怕了，来到\Person{比拉多}面前，恳求他说：\Quote{给我们士兵，让我们看守他的坟墓三天，免得他的门徒来偷走他，百姓以为他从死人中复活了，就对我们行恶。}\Person{比拉多}就把百夫长\Person{沛特洛尼乌斯}连同士兵给了他们，去看守坟墓。长老和经师同他们来到坟墓，与百夫长和士兵一起滚了一块大石头，所有在那里的人一同把它放在坟墓门口；他们贴了七个封条，在那里搭了帐篷守卫。安息日将尽时，清早有一大群人从\Place{耶路撒冷}和周围地区来到，要看那被封住的坟墓。\par

9. 在主的日子即将来临的那个夜晚，当士兵两两站岗时，天上传来巨大的声音；他们看见天开了，有两个人从那里带着大光降下，走近坟墓。那放在门口的石头自己滚开，让出一条路；坟墓打开了，两个青年都进去了。\par

10. 当那些士兵看见后，就叫醒了百夫长和长老们；因为他们也靠近在守卫。当他们讲述所看见的事时，又看见三个人从坟墓里出来，其中两个扶着一个人，一个十字架跟着他们；那两个人的头达到天上，但被他们扶着的那人的头高过诸天。他们听见一个声音从天上说：\Quote{你向那些睡着的人宣讲了。}从十字架上传来回应：\Quote{是的。}\par

11. 他们彼此商议，是否要去把这些事报告给\Person{比拉多}。当他们还在思忖时，天又开了，有一个人降下，进入坟墓。百夫长和与他在一起的人看见这些事，就在夜里急忙赶到\Person{比拉多}那里，离开了他们看守的坟墓，报告了他们所看见的一切事，大为不安地说：\Quote{他真是天主子。}\Person{比拉多}回答说：\Quote{我清白于这天主子的血；但这是你们决定的。}于是他们都上前来恳求他，请求他命令百夫长和士兵不要谈论所看见的事：他们说：\Quote{我们宁可在天主面前犯最大的罪，也不愿落在犹太人的手中被石头砸死。}因此\Person{比拉多}命令百夫长和士兵闭口不言。\par

12. 及至主的日子天将破晓，玛利亚玛达肋纳——主的门徒——因惧怕犹太人，因为他们怒火中烧，尚未在主的坟墓处履行妇女惯常为死者并为所爱之人所行之事——便带着她的朋友同来，到了埋葬他的坟墓那里。她们惧怕被犹太人看见，便说：\Quote{虽然在他被钉的那一天，我们不能哭泣哀悼，但如今让我们在他的坟墓行这些事吧。可是谁为我们滚开墓门前的那块石头，好使我们能进去坐在他旁边，并履行应尽之事呢？因为那石头甚大，我们怕被人看见。倘若我们不能进去，至少我们将带来的纪念物放在门口，哭泣哀悼，直到我们回家为止。}\par

13. 她们前去，发现坟墓已开，靠近往里观看；只见一个少年人坐在坟墓中间，容貌俊美，身穿极其光明的衣袍；他对她们说：\Quote{你们为何而来？你们找谁？是那被钉十字架的吗？他已复活并离开了。倘若你们不信，就进去看看他躺卧的地方，他[不在这里]；因为他已复活，离开这里，往他所被派来的地方去了。}于是妇女们惧怕，便逃跑了。\par

14. 那时正是无酵节的最后一日，许多人正在离去，返回自己的家乡，因为庆节已结束。但我们，主的十二门徒，哭泣忧伤；各人因所发生的事而忧伤，便各自回家去了。但我西满伯多禄和我的兄弟安德肋拿起网，往海边去；同我们一起的还有阿耳斐的儿子肋未，主……\par

\SourceRecord{Translated by J. Armitage Robinson.；Ante-Nicene Fathers；Vol. 9.；Edited by Allan Menzies.；Buffalo, NY: Christian Literature Publishing Co.,；1896.；<http://www.newadvent.org/fathers/1001.htm>.}
